\documentclass[11pt]{article}
\usepackage[a4paper,margin=25mm,headheight=14pt]{geometry}
\usepackage[T1]{fontenc}
\usepackage{lmodern,microtype,amsmath,amssymb,amsthm,mathtools}
\usepackage{booktabs,array,longtable,enumitem,xcolor}
\usepackage{fancyhdr}
\usepackage[colorlinks=true,linkcolor=blue!45!black,urlcolor=blue!50!black,citecolor=blue!45!black]{hyperref}
\usepackage{listings}
\definecolor{ink}{RGB}{26,43,65}
\definecolor{soft}{RGB}{241,245,249}
\hypersetup{pdftitle={K215: verified proof refinements},pdfauthor={Research packet for Vladimir Reshetnikov}}
\pagestyle{fancy}
\fancyhf{}
\fancyhead[L]{\small\textcolor{ink}{K215: verified proof refinements}}
\fancyhead[R]{\small 21 October 2026}
\fancyfoot[C]{\thepage}
\setlength{\parskip}{4pt}
\setlength{\parindent}{0pt}
\setlist{itemsep=2pt,topsep=4pt}
\setcounter{tocdepth}{2}
\lstset{basicstyle=\ttfamily\small,breaklines=true,columns=fullflexible,keepspaces=true,
 backgroundcolor=\color{soft},frame=single,rulecolor=\color{black!15},xleftmargin=3pt,xrightmargin=3pt}
\newcommand{\Con}{\operatorname{Con}}
\newcommand{\Halt}{\operatorname{Halt}}
\newcommand{\NIP}{\mathsf{NIP}}
\newcommand{\DROP}{\mathsf{DROP}}
\newcommand{\ZFC}{\mathsf{ZFC}}
\newcommand{\code}[1]{\texttt{#1}}
\newcommand{\packet}{\texttt{K215\_Proof\_Refinements\_2026-10-21}}
\newtheorem{theorem}{Theorem}[section]
\newtheorem{lemma}[theorem]{Lemma}
\theoremstyle{remark}
\newtheorem{remark}[theorem]{Remark}
\title{\textcolor{ink}{K215: verified proof refinements}\\[5pt]
\large Certificate composition, decoder context propagation,\\
and selection that permits early halting}
\author{Research and verification packet for Vladimir Reshetnikov's ProveIt}
\date{21 October 2026}
\begin{document}
\maketitle
\vspace{-8mm}
\begin{abstract}
The most effective immediate simplification is to retain the accepted 215-state
machine and reduce repeated proof obligations at its interfaces. This packet
implements three such steps. First, a kernel-checked Rocq/Coq adapter combines
the incoming 1,426-entry normalized source certificate with balanced lookup,
preserving the public source-equivalence theorem. Second, a Lean theorem
derives every internal decoder path from one initial marker invariant, removing
separate route and intermediate tape-window premises. Third, an early-halting
selection theorem replaces the nonfalsity requirements used only to ensure
successful extraction of a row from the AX stack.

These are new checked components, not a freshly rebuilt complete K215-to-ZFC
endpoint. The report distinguishes their exact boundaries from the repository's
existing full endpoint. It compares alternative 215-state machines and records
five additional bounded searches, including the final depth-41 R41 table omitted by its original
one-row fold census. No verified 214-state construction was found; the search
failures establish no lower bound.
\end{abstract}

\section{Result and current baseline}

The repository was inspected at
\[
\texttt{aa66d0a33f1f8237a348899eb8095b045a502832}.
\]
This matters because several attractive simplifications were already present
either in maintained sources or in the three newly submitted archives.
The accepted K215 already has Lean and Rocq correctness endpoints. Its
phase-five Lean record reports a 184-module closure, with 37 affected modules
freshly compiled and 147 previously checked objects reused after exact
dependency and source binding. The phase-five work already removes zero
offsets from the Lean source certificate and shares several semantic
arguments \cite{baseline,phase5}. Those changes are background, not claimed
as contributions of this packet.

The three incoming packets at the inspected commit contain further results:
normalized source correspondence, balanced lookup for the older certificate,
arbitrary-tape local decoder blocks, concrete zero-credit enumeration, and
small table-to-table simulations. Their explicit remaining boundaries guided
the work here \cite{checked,refinements,affine}.

\begin{center}
\small
\begin{tabular}{p{34mm}p{65mm}p{38mm}}
\toprule
Component & New result in this packet & Remaining boundary\\
\midrule
Source correspondence &
Balanced lookup composed with the normalized 1,426-entry checker; exact
all-register theorem and public facade compiled &
Production Rocq 9.1 endpoint rebuild\\
\addlinespace
Decoder &
All internal source paths follow from one marker and actual pointer bits;
no per-edge window proofs &
Maintained-definition identity, primitive edges, returns\\
\addlinespace
AX selection &
A generic row-selection theorem permits early halting and preserves the
counter and opaque stack tail &
Integration scope detailed in Section~\ref{sec:selection}\\
\addlinespace
State reduction &
Five bounded searches on exact current tables, with complete receipts &
No 214-state construction; no lower bound\\
\bottomrule
\end{tabular}
\end{center}

The recommendation is to keep primary K215 as the main formal endpoint.
For a second 215-state theorem, a small marker-transfer alternative is the
least demanding incremental proof. A zero-credit machine has a cleaner
independent enumeration argument, but changing the scheduler creates a
larger integration task than retaining the already verified program.

\subsection{What the checked statements mean}

The desired mathematical statement remains
\[
\ZFC\vdash
\bigl(\neg\Halt(K215,\mathrm{blank})\leftrightarrow\Con(\ZFC)\bigr).
\]
It identifies nonhalting with syntactic consistency; it does not assert either
side unconditionally. The present proof refinements preserve the source
programs or the relevant execution contracts. They introduce no new
set-theoretic assumption.

A theorem about a source program, a finite decoder prefix, or a command
macro is not by itself a complete binary-machine endpoint. The compiled
results below state which interface they cover. Likewise, a checked external
Lean or Rocq theorem is not a new internal PA derivation. The repository's
separate PA status is unchanged.

\section{Simplify the interfaces before changing the calculus}

The current endpoint factors through four different kinds of reasoning:
\[
\begin{aligned}
\Halt(K215,\mathrm{blank})
&\longleftrightarrow \Halt(p_{215},R_0),\\
\Halt(p_{215},R_0)
&\longleftrightarrow \Halt(\mathrm{dagPerm},R_0),\\
\Halt(\mathrm{dagPerm},R_0)
&\longleftrightarrow \text{a finite refutation in the calibrated calculus},\\
\text{a finite calculus refutation}
&\longleftrightarrow \text{a finite ZFC refutation}.
\end{aligned}
\]
The first line contains the physical tape, initialization, and compiler
simulation. The second is source correspondence. The third is the
enumeration and proof-compilation argument. The last is logical calibration.

These layers should communicate through compact theorems. Repeating a
source simulation, a tape-context reconstruction, or a list of syntactic
nonfalsity facts in each client increases the number of obligations without
strengthening the endpoint.

The guiding method in this packet is \emph{checked refinement}. A new
representation or finite checker is related to an existing sound interface
by a small theorem. An untrusted generator may compute a tree, role table,
or certificate; the kernel checks the finite conditions on its actual
output. Correctness does not depend on trusting the generator.

This approach also makes future machine changes cheaper to assess. A new
table can reuse source and logical calibration if it satisfies the same
execution contracts. Conversely, a new calculus must supply its own
refutation translations; a shorter source listing alone does not establish
equiconsistency.

\section{Balanced lookup for the normalized source certificate}
\label{sec:source}

\subsection{The two existing refinements}

The incoming normalized checker identifies instruction pointers that differ
only by trailing zero bits. With zero padding understood,
\[
P\approx Q
\quad\Longleftrightarrow\quad
\forall i\in\mathbb N,\ \operatorname{bit}(P,i)=\operatorname{bit}(Q,i).
\]
This removes duplicate checkpoint representations, reducing the ordinary
certificate from 2,240 entries to 1,426. The incoming proof transports
completed segment endpoints through the established pointer-equivalence
theorem. It does not blindly replace pointer arithmetic by arithmetic on
trimmed lists.

Between two guarded decrement checkpoints, a finite segment accumulates
increments. The checker compares the resulting register-increment counts
and the next normalized pointers. At a decrement checkpoint it checks the
same register and both branches. The soundness theorem quantifies over every
natural-valued register file, not only reachable small values.

The other incoming refinement stores the older 2,240-entry certificate in a
balanced binary tree. Its benefit is lookup work: an index is resolved in a
number of tree nodes proportional to the tree height, rather than by walking
a linked-list prefix. The earlier packet explicitly left composition with
the normalized checker as further work \cite{refinements}.

\subsection{A small refinement theorem composes them}

A certificate tree has empty, leaf, and internal-node constructors. An
internal node stores a binary natural number giving the size of its left
subtree. Let \(\operatorname{flat}(T)\) be the in-order list of entries.
A validity predicate checks that every stored cut equals the actual left
subtree size.

The inherited tree library proves
\[
\operatorname{lookup}(T,j)
=
\operatorname{nth\_error}\bigl(\operatorname{flat}(T),\operatorname{toNat}(j)\bigr)
\qquad\text{when }T\text{ is valid}.
\]
Out-of-range indices are included: both sides return failure. No default
certificate entry is silently accepted.

The new \code{BalancedNormalized.v} uses tree lookup inside the normalized
branch and entry checkers. It proves the following exact Boolean equality.

\begin{theorem}[Balanced normalized refinement]
For every valid certificate tree \(T\) and source trees \(t_1,t_2\),
\[
\operatorname{nbcheck}(t_1,t_2,T)
=
\operatorname{ncheck}\bigl(t_1,t_2,\operatorname{flat}(T)\bigr).
\]
Consequently, successful tree validity and successful balanced normalized
checking imply the existing all-register source halting equivalence.
\end{theorem}

\begin{proof}
Replace tree lookup by list lookup in each branch. The two normalized
pointer comparisons and the register-count comparison are unchanged.
The entry equality follows by inspecting the two source instructions;
the decrement case retains both branches. Tree-wide conjunction equals
list-wide conjunction on the flattened entries. Apply the existing
one-sided normalized checker theorem, whose mirror argument supplies the
reverse implication.
\end{proof}

The adapter is short because it adds no new operational semantics. The
unbounded simulation proof remains the already checked normalized one.

\subsection{The exact K215 instance}

The concrete tree is generated from the existing \code{Normalized215}
literal. Separate kernel-checked statements establish:
\[
\begin{gathered}
\operatorname{flat}(T_{215})=\operatorname{Normalized215},\\
\operatorname{validb}(T_{215})=\mathrm{true},\qquad
\operatorname{size}(T_{215})=1426,\qquad
\operatorname{height}(T_{215})=12.
\end{gathered}
\]
Every lookup has at most 12 visited tree nodes. The theorem about height
includes arbitrary indices; the actual checker additionally rejects failed
lookups.

The resulting exact source theorem is
\[
\forall R:\mathbb N\to\mathbb N,\quad
\operatorname{rm\_halts}(p_{215},([],R))
\longleftrightarrow
\operatorname{rm\_halts}(\mathrm{dagPerm},([],R)).
\]
The facade exports the existing name
\code{K215F.Source215.p215\_dagPerm}. An explicit typed caller compiled
against that facade. The maintained \code{Correct215.v} uses precisely
this interface when composing source and logical correctness
\cite{sourceendpoint}.

Both literal source files, \code{Prog215.v} and \code{DagPerm.v}, were fetched
at the inspected commit and compared byte for byte with the compiled inputs.
The package retains their hashes and the current public boundary in
\code{balanced\_normalized/current-binding.json}.

\subsection{Exact lookup measurements}

The following census counts only the data-structure nodes visited by
certificate endpoint lookup. It does not count source-tree decoding,
segment evaluation, proof elaboration, tree construction, or kernel work
outside this operation.

\begin{center}
\small
\begin{tabular}{lrrrr}
\toprule
Representation & Entries & Lookups & Visited nodes & Maximum tree nodes\\
\midrule
Ordinary list & 2,240 & 4,479 & 5,018,706 & --\\
Ordinary balanced tree & 2,240 & 4,479 & 54,514 & 13\\
Normalized list & 1,426 & 2,851 & 2,033,059 & --\\
Normalized balanced tree & 1,426 & 2,851 & 32,961 & 12\\
\bottomrule
\end{tabular}
\end{center}

There are 1,425 decrement entries and one halt entry in the normalized
certificate. Thus the root check plus two references per decrement gives
\(1+2(1425)=2851\) lookups.

The normalized data file is 108,633 bytes, compared with 171,678 bytes for
the ordinary compact data file. The shared path codec adds 1,325 bytes.
These are certificate-source measurements, not total proof-library sizes.
Balanced storage reduces a cost inside checking; it does not by itself
reduce the number of machine states or prove a whole-build speedup.

\subsection{Fresh verification and scope}

The isolated package freshly compiled 17 modules with Coq 8.18.0 and OCaml
4.14.1. The nine assumption reports printed by the new adapter, concrete
tree statements, exact theorem, facade, and client are closed under the
global context. The summed compiler time in this execution was 50.941
seconds; the exact \code{Source215Composed.v} compilation took 9.126 seconds.
These are observations of this run, not a benchmark comparison.

A separate recursive
\code{coqchk -silent EndpointClient} completed successfully in 306.215
seconds. It checked the concrete client and its recursive dependencies,
rather than only the abstract tree library.

The compatibility route adds the supplied \code{Compat818} import to the
inherited \code{RMEquiv.v} source. That exact compiled source is hashed in
the receipt. No new axiom is introduced. A maintained Rocq 9.1 build and
the full binary/ZFC endpoint were not rebuilt here. The package includes
the production-toolchain branch of the reproduction recipe, but it is not
represented as executed evidence.

\section{Derive decoder paths from one invariant}
\label{sec:decoder}

\subsection{What the local-block proof still required}

The incoming decoder package already proves 281 local finite-window blocks
for arbitrary complete tapes. It also proves a chain-composition theorem.
However, a client of that chain theorem still had to supply the semantic
window premise for each successive block. Enumerating the blocks does not
itself prove that their required tape contexts arise on a real source path.

Two modules close this obligation.
The generic proof is in \code{DecoderContext.lean}.
Its concrete instance is in \code{K215Context.lean}.
They use the three unchanged incoming foundations. A typed audit module
makes six fresh compilations in the delivered check.

\subsection{Finite contexts and their propagation}

A decoder role consists of a source node, a physical state, a finite suffix
immediately preceding the head, and its two outgoing edges. Different
roles may use the same physical state in different contexts.

Write \(\operatorname{Before}(s)\) for the finite partial tape that places
a suffix \(s\) at positions
\(-|s|,\ldots,-1\). For the current input bit \(b\), the known local context is
\[
K(r,b)=\operatorname{Before}(\operatorname{suffix}(r))\cup\{0\mapsto b\}.
\]
A finite coverage check confirms that every cell required by a block's
input window occurs in \(K(r,b)\) with the same value. A second coverage
check confirms that \(K(r,b)\) supplies the next role's suffix after the
head moves one cell to the right. Entry and exit states are checked too.

These conditions compose because the inherited local theorem restores
the \emph{entire} tape. Intermediate head movements and writes are allowed;
their net effect is the certified state change and the one-cell head
advance. Consequently the next suffix can be proved from the same original
tape.

\begin{lemma}[One internal edge]
Suppose the local block and its finite context checks pass.
If the complete tape matches the current suffix and current bit at head
position \(a\), the block runs for positive time, ends in the next role at
\(a+1\), restores the whole tape, and establishes the next suffix invariant.
\end{lemma}

\begin{proof}
Finite cell inclusion supplies the local block's arbitrary-tape window
premise. Apply the inherited block theorem. Since its output tape is
identical to its input tape, the second finite inclusion supplies the
shifted suffix. Positivity comes from the checked block duration.
\end{proof}

The checker verifies these conditions for all 371 internal edges in the
234-role product, using 742 context inclusions. Primitive edges are handled
separately and are deliberately outside this check's semantic claim.

\subsection{A source-path theorem without a route witness}

A further finite check links the role graph to the supplied 202-node source
DAG. Whenever a source child is internal, the corresponding role edge must
be internal and point to a role for that child. This proves that every
internal source path has a decoder route. It removes even an explicit
route or successful-follow premise from the final theorem.

Let \(\operatorname{Reads}(T,a,w)\) mean that the successive tape cells
starting at \(a\) contain the bits of \(w\).
The exact root role has physical state 84 and suffix \([1]\).

\begin{theorem}[Every internal K215 source path]
\label{thm:decoder}
Suppose \(w\) follows the supplied source DAG from its root to an internal
node \(m\). For every integer \(a\) and complete Boolean tape \(T\), assume
\[
T(a-1)=1,\qquad \operatorname{Reads}(T,a,w).
\]
There exist a role \(f\) and a duration \(n\) such that \(f\)'s source node is
\(m\), and execution of the supplied literal K215 table satisfies
\[
(84,a,T)\ \xrightarrow{\ n\ }\ (f.\mathrm{state},a+|w|,T).
\]
The next-role suffix invariant holds, and \(w\ne[]\) implies \(n>0\).
\end{theorem}

\begin{proof}
Induct on the source path, using the finite source-link check to obtain
the successive role edges. Apply the one-edge lemma at each step and
compose finite executions. Tape restoration supplies all subsequent
pointer bits and suffix facts. Nonempty paths contain a positive-duration
block. The empty path has duration zero.
\end{proof}

The compiled theorem is
\code{K215Context.every\_internal\_source\_path}.
Its tape quantifier is unrestricted outside the initial marker and supplied
pointer bits. It does not assume a finite-support tape, a register bound,
or a successful finite simulation of some particular run.

\subsection{Concrete checking and remaining adapter}

The six-module Lean 4.32.0 build completed with zero warnings and only the
standard Lean axioms in the printed semantic theorem reports.
No \code{sorryAx} or \code{native\_decide} is used. Five checked malformed
controls reject a missing predecessor fact, a wrong predecessor fact, a
wrong next suffix, a missing target role, and a primitive source node used
as an internal parent.

An independent finite audit binds every published transition row and every
expanded source leaf to the supplied constants: 430 rows and 4,751 leaves.
It also checks the 4,750 internal source occurrences, including the root.
The longest internal prefix has 46 bits and takes 55 steps in this exact
decoder. These finite measurements corroborate the literal instance; the
universal path statement is the kernel theorem.

This is a substantial strengthening of the incoming chain interface, but
it is not yet a replacement for the whole runtime proof. Three tasks remain:
\begin{enumerate}
\item Prove in the maintained kernel namespace that the supplied native table
and source-DAG constants denote the maintained K215 definitions.
The delivered external equality audit is exact, but is not that kernel
adapter. Transport the complete integer-indexed tape semantics to the
maintained tape representation as part of this binding.
\item Connect the final primitive-entry edges to the existing primitive
register-operation contracts. Theorem~\ref{thm:decoder} ends at an internal
source node.
\item Reuse the return and register-frame theorems to reestablish the root
marker, and retain the exact blank initialization proof.
\end{enumerate}
The value of the new theorem is that none of these clients needs to rebuild
intermediate decoder window proofs for every occurrence.


\section{Selection may halt early}
\label{sec:selection}

\subsection{The unnecessary strength in the old selection argument}

A completeness compiler needs a block that either reaches its requested
formula configuration or halts earlier. Early detection of the falsity code
already achieves its purpose. Proving that every intermediate axiom or gate
code is different from falsity is therefore unnecessary when those facts
are used only to force normal return from a selector.

This distinction is especially useful for the AX bundle: selecting one row
should be a generic stack operation. It should not require a separate
syntactic proof about every skipped row. The theorem below allows arbitrary
row values, including non-formula codes and the falsity code itself.

\subsection{The exact command semantics}

Write a numeric machine state as \((T,W,D,H)\), representing top, stack,
counter, and cache. Let \(P(x,W)\) be stack push, and let
\(\operatorname{unP}\) be its inverse, with the established law
\[
\operatorname{unP}(P(x,W))=(x,W).
\]
The literal production command shapes are \cite{compiler}
\[
\NIP(T,W,D,H)=
\begin{cases}
\mathrm{halt},&H=1,\\
(T,\pi_2(W),D,\pi_1(W)),&H\ne1,
\end{cases}
\]
and
\[
\DROP(T,W,D,H)=(H,P(T,W),D,H).
\]
In particular, NIP tests the \emph{old cache before popping}. This timing
determines the correct selection counts.

Let \(\operatorname{pack}(L,W)\) push the explicit list \(L\) above the
opaque numeric tail \(W\), with the first listed value popped first.

\begin{samepage}
\begin{theorem}[Early-halting indexed selection]
\label{thm:selection}
Let \(L=[x_0,\ldots,x_{r-1}]\) and \(0\le j<r\).
From state
\(\bigl(G,\operatorname{pack}(L\mathbin{+\!\!+}[0],W),D,H\bigr)\),
the command word
\[
\NIP^{\,j+1}\ ;\ \DROP\ ;\ \NIP^{\,r+1-j}
\]
either halts early or returns exactly \((x_j,W,D,0)\).
There are no theoremhood, formula-well-formedness, nonfalsity, or
fresh-variable premises on \(G,H\), or the entries of \(L\).
\end{theorem}
\end{samepage}

\begin{proof}
If any alarm occurs, the first alternative holds. Otherwise, the first
\(j+1\) NIPs leave \(x_j\) in the cache and remove the preceding rows and
\(x_j\) from the explicit stack prefix. DROP makes \(x_j\) the top, retains
it in the cache, and pushes the old top \(G\).

The remaining explicit prefix now consists of \(G\), the \(r-1-j\)
remaining rows, and the final zero: exactly \(r+1-j\) values. If no alarm
occurs during these NIPs, their last pop places zero in the cache. The top
remains \(x_j\), the counter remains \(D\), and the stack is exactly \(W\).
The tail \(W\) is never popped. The final zero is loaded but is not tested
again within this block.
\end{proof}

The generic proof uses only the push/pop inverse law. It does not rely
on a separately stored stack length, so numerical aliases such as
\(P(0,0)=0\) cause no problem. The explicit-prefix calculation is justified
by repeated applications of that inverse law.

The extra zero is essential to the zero-cache conclusion. Replacing it by
an arbitrary nonzero sentinel changes the final cache on a normal run.
The finite corroboration includes a counterexample to that tempting
overgeneralization.

\subsection{Exact binding to the maintained command language}

This packet goes beyond a second, separately defined command semantics.
The compiled \code{K215SelectorAdapter.v} proves:
\begin{itemize}
\item the explicit embedding of NIP and DROP into the production
\code{L1\_CompPerm.cmd};
\item exact one-step and execution equalities for that embedding;
\item the concrete selection-word equality;
\item the exact nine-row view of production \code{EnumPerm.ax\_stack};
\item the actual-command theorems \code{select\_actual} and
\code{select\_AX\_row}.
\end{itemize}

The nine rows, in pop order, are
\[
\begin{aligned}
[&\rho_1(a,b),\rho_2(a,b,c),\rho_3(a,b,c),\rho_4(a,b,c),\\
 &\rho_5(a,b,c),\rho_6(a,c),\operatorname{growG}(a,b,c),
 \rho_8(a,b,c,u),\operatorname{wal}(u,h)].
\end{aligned}
\]
They are followed by the saved old top. After the existing reset/discard
setup, that saved top is zero. The eight axiom selectors use indices
\(0,\ldots,7\); GEN uses index 8 after its existing setup establishes the
same zero-bottom condition. For these nine rows every selector has
\(9+3=12\) commands, although the number of NIPs on each side of DROP varies.

The production theorem states that, from the actual AX stack with saved
top zero, production \code{exec} either returns \code{None} or returns the
selected row with the original tail and zero counter/cache. No definition
of the production commands or AX stack is changed.

\subsection{Verification and how to adopt the result}

The two generic modules compiled and passed recursive \code{coqchk}.
A second fresh build compiled all eleven modules in the production adapter's
small dependency closure, including the current command compiler and
enumeration definitions, and recursively rechecked
\code{K215SelectorAdapter}. The two new actual-command assumption reports
are closed.

Five current production source files were bound to their GitHub blobs at
the inspected commit. For the Coq 8.18 compatibility run, three receive
only an import-prefix adjustment from \code{Stdlib} to \code{Coq}.
The inherited optimized dominated-rho8 proof is explicitly credited and
freshly compiled in this dependency closure. The earlier failed adapter
attempt, due to a missing explicit import, is retained separately from the
successful final receipt.

The independent numeric check covers 11,502 finite cases: 10,326 early
halts and 1,176 exact returns. It checks zero aliases, small arbitrary tails,
counter preservation, and the sentinel counterexample. These cases are
corroboration; the universal theorem and its actual-command binding are
kernel-checked.

The next production edit is small in scope: retain the exact discard/setup
contracts and replace the eight row-extraction arguments and GEN extraction
tail with \code{select\_AX\_row}. The declared completeness contract already
allows early halting. Keep nonfalsity lemmas used by exact-return clients
or elsewhere in soundness.

The current leaf proofs already use a compact tactic, so a reduction in
total source lines or kernel time has not been established. What is
proved is that the stronger nonfalsity obligations are unnecessary on
this selection path. The full maintained \code{build} theorem and final
binary/ZFC endpoint have not been replaced or rebuilt.


\section{Which other 215-state candidate is easier to prove?}
\label{sec:alternatives}

There are three different meanings of \emph{easier} here: smallest additional
proof over the accepted endpoint, cleanest independent enumerator, and
simplest physical encoding. They point to different candidates.

\begin{center}
\small
\begin{tabular}{p{24mm}p{52mm}p{62mm}}
\toprule
Candidate & Best reason to use it & Qualification\\
\midrule
Primary K215 &
Already has the complete formal endpoint; all three refinements above
apply without changing its table &
Best choice for reducing present maintenance work\\
\addlinespace
M124 / FM215 &
Small marker-based table simulation transports the existing correctness
argument &
Representation bridge and endpoint composition still need their own checks\\
\addlinespace
CM215 &
Incoming work already records a Lean Core bridge &
Least remaining Lean adapter work may outweigh a more involved local invariant\\
\addlinespace
ZC215 &
Zero-credit enumeration removes the reserve/debt calculation &
Its physical runtime and changed scheduler still need full integration\\
\addlinespace
R41 &
Shallower source addresses, same calculus and primitive control flow &
Shorter addresses are not a smaller state count or a general speedup\\
\bottomrule
\end{tabular}
\end{center}

\subsection{The least additional proof: a marker-transfer machine}

For a second theorem derived from primary K215, M124 or FM215 is a good
choice. The incoming finite-table theorem already isolates two predecessor
marker facts and a one-or-three-step simulation on arbitrary tapes
\cite{refinements}. It avoids proving a new calculus or scheduler.

The concrete family uses old roles 124 and 127. Incoming transitions
establish predecessor markers 1 and 0, respectively. A shared role and a
selector distinguish the two old zero branches by that marker, returning
through existing provider states. The providers differ between the
literal M124 and FM215 tables. The local simulation preserves the complete
tape and head at each corresponding boundary and uses positive time.

The invariant concerns the immediately preceding cell at two control
states, so it is far smaller than a global register-format invariant.
Since the roles are noninitial and all their incoming transitions write
the required marker, the relation starts from arbitrary tapes at state zero.

The alternative remains at 215 states because one role is merged but
the selector occupies the released role's state. The operation reorganizes
the decoder; it does not yield a net deletion. Reporting it as a 214-state
construction would miscount the selector.

These are existing alternatives, not new machines discovered in this
packet. The incoming native table-to-table theorems do not by themselves
constitute the repository's dense/zipper ZFC endpoint. CM215 may be faster
to integrate in Lean because an incoming packet already records its Core
adapter. The exact next endpoint should be chosen with that engineering
status in mind.

\subsection{The cleanest independent scheduler: ZC215}

ZC215 assigns DROP the lower tag and NIP the next tag. Lowering a marked
NIP occurrence to the corresponding DROP occurrence gives an injection into DROP occurrences belonging to earlier raw
codes in the exhaustive schedule. Thus earlier DROP
credits cover the potential NIP debits. The incoming packets now prove
the concrete Cantor-code and decoder facts behind this argument, rather
than leaving them abstract \cite{checked,refinements}.

This is a useful simplification of an independently reconstructed
enumerator. It does not automatically minimize the total proof needed
today, because primary K215 already has a complete endpoint. The combined
GC-ZC alternative also changes the binder/compiler arrangement and hence
introduces more interfaces to revisit.

The new early-halting selector is compatible with this general proof
strategy, but it does not alter command tags or establish a new binary
machine. A proof simplification and a state reduction remain different
claims.

\subsection{The encoding comparison: R41}

The R41 source retains 202 distinct subtrees and 4,751 primitive occurrences.
Its maximum address depth is 41 instead of 47, and the sum of address lengths
is 149,265 instead of 160,305. Nineteen transition rows change, and the
published source/runtime evidence remains at 215 working states
\cite{r41}.

The reduction of 11,040 address bits is a static improvement. The original
R41 packet explicitly reports mixed runtime changes, so it should not
be advertised as uniformly faster. Nor does a shallower tree make its
correctness proof as small as a local transfer from the already checked
primary table. It is nevertheless a useful source for testing a different
allocation geometry, which motivated the new searches below.

\section{The search for 214 states}
\label{sec:search}

No verified 214-state machine was constructed in this work. A useful
coverage gap was found in the original R41 one-row fold census: it tested
32 earlier depth-44/45 tables and explicitly excluded the final depth-41
machine \cite{r41}. This packet tests the actual final table in that family,
along with two newer layouts. Other incoming investigations have tested
R41 under different families; no claim of first-ever R41 search is made.

\begin{center}
\small
\begin{tabular}{p{48mm}p{56mm}r}
\toprule
Exact table and family & Bounds & Specifications\\
\midrule
R41, one-row contextual &
History depth 4; 45 steps; 200 symbolic nodes & 71,380\\
R41, one-row contextual &
History depth 8; 97 steps; 1,000 symbolic nodes & 99,760\\
Shallow context-return215, one-row contextual &
Depth 8; 97 steps; 1,000 symbolic nodes & 108,360\\
GC-ZC215 canonical, one-row contextual &
Depth 8; 97 steps; 1,000 symbolic nodes & 97,180\\
R41, blank-window folds &
Radius 7; 97 steps; specified inherited/common-action family & 632,584\\
\bottomrule
\end{tabular}
\end{center}

All five searches returned no survivor. The counts describe search
specifications, not distinct complete machines; overlapping families
should not be added and advertised as a count of unique constructions.

The first four families require the opposite-symbol transitions of each
candidate state pair to be literally equal. They enumerate replacement
writes, moves, and working-state targets and use bounded graph histories
as local contexts.

The final search exhausts all 46,010 ordered pairs of distinct working
states, including initial-state mergers. It considers 182,992 inherited
action choices and 449,592 one-common-action choices. All 632,584 fail its
finite-window common-endpoint filter; none reaches the independent
acceptance checker. The observed execution finished before its inherited
110-second ceiling, so that ceiling did not truncate this recorded run.
The complete compressed rejection record is included.

The three input tables and both main search entry points were fetched at
the inspected GitHub commit and matched byte for byte to the executed
copies. The source algorithms are inherited unchanged; the new contribution
is the explicitly identified table coverage and preserved executions.

\subsection{What these failures do not establish}

Failure of a bounded local filter is not a proof of semantic inequivalence.
A candidate may require a longer reconciliation path or a sharper invariant
than the filter supplies. A history overapproximation may include impossible
contexts and reject a sound machine. Larger coordinated edits, another
source layout, and a different calibrated calculus remain outside these
searches.

The results also do not replace the incoming packet's much larger
coordinated two-merge/one-selector investigation. The families have different
restrictions. None supplies a lower bound on the state count.

\subsection{A concrete next search direction}

The new decoder theorem suggests a better proof interface for future
allocation experiments. Search can operate on roles together with finite
suffix contexts. A successful candidate should export:
\begin{enumerate}
\item a complete literal transition table, including blank initialization;
\item the source-node/physical-state/context role product;
\item positive-duration blocks with whole-tape effects;
\item finite context-inclusion and source-link checks;
\item primitive and return contracts plus the exact initial frame.
\end{enumerate}
The kernel can then validate the candidate's actual contexts instead of
requiring a fresh expanded proof for every source occurrence.

To get a net state saving from the two-marker family, the discriminator
must be supplied without consuming the freed state, or a coordinated
change must release another state elsewhere. The obvious short local
versions are heavily searched. A useful new attempt should jointly vary
source layout, physical allocation, and the context invariant, and should
count the selector, providers, and initialization before claiming 214.

On the formalism side, the permission to generate junk codes is valuable
only when the altered proof system still has explicit refutation
translations in both directions. The selection theorem shows one place
where junk values need no extra syntactic case analysis: if they trigger
the designated halt, completeness has already succeeded. It does not
license unsound inference rules or supply a new equiconsistency proof.
No smaller calibrated calculus is claimed here.

\section{Integration plan and exact acceptance gates}

The new results are useful as a sequence of small changes.

\paragraph{1. Adopt the source facade first.}
Move the combined checker and data into the maintained Rocq source
namespace, update the generator and dependency inventory, and preserve
\code{K215F.Source215.p215\_dagPerm}. Rebuild the production
\code{Correct215} and \code{CorrectDense215} modules under the maintained
toolchain and inspect their assumptions. The checked facade and client
make this a concrete integration task; they do not replace that final
production build.

\paragraph{2. Use the actual-command selector.}
Replace the eight leaf-extraction arguments and the GEN extraction tail
with the compiled \code{select\_AX\_row} theorem. Retain the existing setup
and exact discard contracts. Recheck the compiler's existing completeness
statement and soundness dependencies. Measure total source/proof-object
size or compile time only after that adoption if a performance claim is
wanted.

\paragraph{3. Install the decoder adapter.}
Identify the literal DAG/table with maintained constants in the kernel
and transport the full-tape semantics into the production tape model.
Compose the all-internal-path theorem with primitive, return, and register
contracts. Keep initialization separate. This is where substantial
occurrence-by-occurrence context proof work can be removed, but the
delivered theorem does not yet quantify that total reduction.

\paragraph{4. Prove one alternative by transport.}
For a second 215-state endpoint, use M124/FM215's local marker theorem
or CM215's already advanced Lean bridge. Avoid rebuilding logical
calibration when a table simulation preserves the existing source
program's behavior. A clean ZC215 endpoint remains worthwhile as an
independent scheduler implementation.

\paragraph{5. Accept a state record only at the full interface.}
A 214-state table must include initialization, have the exact working-state
count, and satisfy an unbounded correctness argument for its own literal
rows and initial tape. Finite startups, static DAG counts, or a bounded
search survivor are not substitutes.

\section{Contents and reproduction}

The archive is organized as follows.

\begin{center}
\small
\begin{tabular}{p{48mm}p{91mm}}
\toprule
Path & Purpose\\
\midrule
\code{report.tex}, \code{report.pdf} &
Editable report and its compiled, visually checked PDF\\
\code{balanced\_normalized/} &
New composition, inherited dependencies, exact facade/client,
metrics, current-source binding, compiler and recursive-audit evidence\\
\code{decoder\_context/} &
New generic and literal Lean context proofs, inherited foundations,
literal inputs, typed audit, and complete fresh-build evidence\\
\code{early\_halt/} &
Generic selector, numeric layer, exact production adapter, source
provenance, fresh builds, and finite corroboration\\
\code{search214/} &
Five bounded search runs, exact input tables, inherited algorithms,
portable replay, and complete compressed rejection records\\
\code{reviews/} &
Scoped independent reviews within this research collaboration\\
\code{MANIFEST.json}, \code{verify\_packet.py} &
Delivered-file integrity and recorded-evidence checks\\
\bottomrule
\end{tabular}
\end{center}

The lightweight audit checks package integrity, recomputes lookup metrics,
checks the recorded proof/source bindings, and reruns the small context and
numeric corroborations. It does not silently invoke a full theorem-prover
build or repeat the large searches.

\begin{lstlisting}
python3 verify_packet.py
\end{lstlisting}

Fresh formal checks can be run separately, using empty output directories:

\begin{lstlisting}
python3 balanced_normalized/build.py \
  --out /tmp/k215-balanced-normalized-fresh --recursive

python3 decoder_context/build.py \
  --lean /absolute/path/to/lean --out /tmp/k215-context-fresh

python3 early_halt/build.py --output /tmp/k215-selector-fresh
python3 early_halt/build_production.py \
  --output /tmp/k215-selector-production-fresh
\end{lstlisting}

The source package accepts Coq 8.18 and includes a Rocq 9.1 recipe branch;
the executed evidence is from Coq 8.18. The decoder requires Lean 4.32.0.
Toolchain installations and compiled proof objects are not included.
The proof sources and the inherited dependency sources needed for these
focused checks are included, along with their notices.

The optional search replay is:

\begin{lstlisting}
python3 -B search214/reproduce.py --output /tmp/k215-fold-replay
\end{lstlisting}

Its blank-window component retains the inherited time ceiling. On a slower
machine the driver may report that the family did not complete, rather
than promoting an incomplete execution to the recorded completion claim.

Build the PDF with the included \code{build\_report.py}, which invokes
pdfLaTeX twice in a separate build directory. Historical logs retain their
actual execution paths as provenance; reproduction entry points use
package-relative inputs.

\section*{Attribution and result boundary}

The accepted calculus, K215 machine, register compiler, earlier proof
certificates, decoder blocks, alternative tables, and search algorithms
retain their original authorship and licenses. The new contributions are
the balanced/normalized composition, automatic internal decoder-context
propagation, the early-halting selection theorem and exact production
adapter, and the new explicitly scoped search executions. The three
incoming packets were inspected before these extensions were developed.

Separate review here means another worker in this research collaboration,
not an independent external human endorsement. Every assertion of formal
checking is tied to the included source hashes, commands, return codes,
and logs.

The strongest new result is a smaller set of reusable proof obligations
for the existing machine. The main K215 endpoint remains inherited.
There is no claimed 214-state machine, unconditional nonhalting theorem,
or minimum-state theorem.

\begin{thebibliography}{9}
\small\raggedright
\bibitem{baseline}
ProveIt, \emph{Smaller machines for Con(ZFC)}, inspected commit
\code{aa66d0a33f1f8237a348899eb8095b045a502832}.
\url{https://github.com/VladimirReshetnikov/ProveIt/blob/aa66d0a33f1f8237a348899eb8095b045a502832/Computability/BusyBeaver/Research/ZFC/README.md}

\bibitem{phase5}
ProveIt, \emph{Ordinary certificates and shared semantic lemmas},
phase-five checkpoint.
\url{https://github.com/VladimirReshetnikov/ProveIt/blob/aa66d0a33f1f8237a348899eb8095b045a502832/Computability/BusyBeaver/Research/ZFC/k225-formal/verification/k215/simplification/phase5/README.md}

\bibitem{checked}
\emph{Checked proof simplifications for K215},
incoming research packet, 21 October 2026. Source normalization,
local decoder certificates, concrete zero-credit argument, and bounded
coordinated state search.
\url{https://github.com/VladimirReshetnikov/ProveIt/blob/aa66d0a33f1f8237a348899eb8095b045a502832/docs/incoming/K215_Checked_Simplifications_2026-10-21.zip}

\bibitem{refinements}
\emph{Verified certificate refinements for K215},
incoming research packet, 21 October 2026. Balanced lookup, shared marker
certificates, concrete credit theorem, and restricted fold searches.
\url{https://github.com/VladimirReshetnikov/ProveIt/blob/aa66d0a33f1f8237a348899eb8095b045a502832/docs/incoming/K215_Certificate_Refinements_2026-10-21.zip}

\bibitem{affine}
\emph{Affine proof cores and incremental machine transport},
incoming research packet, 21 October 2026.
\url{https://github.com/VladimirReshetnikov/ProveIt/blob/aa66d0a33f1f8237a348899eb8095b045a502832/docs/incoming/Affine_Proof_Cores_and_Transport_2026-10-21.zip}

\bibitem{sourceendpoint}
ProveIt, \code{K215F.Correct215}, maintained source-to-logical endpoint.
\url{https://github.com/VladimirReshetnikov/ProveIt/blob/aa66d0a33f1f8237a348899eb8095b045a502832/Computability/BusyBeaver/Research/ZFC/k225-formal/rocq/rocq-extension/coq/215/Correct215.v}

\bibitem{r41}
ProveIt, \emph{K215-R41: a rebalanced 215-state ZFC-consistency machine}.
\url{https://github.com/VladimirReshetnikov/ProveIt/blob/aa66d0a33f1f8237a348899eb8095b045a502832/Computability/BusyBeaver/Research/ZFC/candidate-k215/alternatives/rebalanced/README.md}

\bibitem{compiler}
ProveIt, \code{L1\_CompPerm.v}, maintained command compiler.
\url{https://github.com/VladimirReshetnikov/ProveIt/blob/aa66d0a33f1f8237a348899eb8095b045a502832/Computability/BusyBeaver/Research/ZFC/k225-formal/rocqperm/L1_CompPerm.v}
\end{thebibliography}
\end{document}

